% ============================================================
% Chapter 1: Introduction
% ============================================================

\chapter{Introduction}
\label{chap:intro}

% ------------------------------------------------------------
\section{Background and Motivation}
\label{sec:intro:background}
% ------------------------------------------------------------

Software development has changed rapidly in recent years. Since GitHub Copilot
launched in October 2021, AI tools have moved well beyond simple code
suggestions. Today's coding agents can plan tasks, write and run code, execute
tests, and open pull requests with little human input. Tools such as
Anthropic's Claude Code, Cursor, Windsurf, OpenAI Codex CLI, and Amazon Q
Developer are now actively used in both open-source and commercial projects.

The pace of adoption is striking. GitHub's Octoverse 2025 report recorded
TypeScript as the most-used language on the platform for the first time, with
the JavaScript and TypeScript ecosystem making up the single largest share of
public activity~\cite{octoverse2025}. This makes the \emph{npm}
ecosystem---the package registry behind both languages---a natural place to
study how AI coding tools are being adopted. npm repositories leave behind
observable traces: lockfiles, dependency declarations, configuration files,
and commit histories that can all be analysed at scale.

Despite how fast these tools are spreading, we still do not have a clear
picture of how many developers actually use them or who those developers are.
Robbes and colleagues produced the first large-scale estimates---22--28\%
adoption globally in established projects, rising to 71--76\% in newly
created repositories~\cite{robbesAgenticMuch2026, robbesAgenticVery2026}---but
their studies cover all GitHub languages and do not focus specifically on
JS/TS or npm. Khosravani and Mockus~\cite{khosravani2026} found that earlier
studies using only one detection signal may miss up to 30 times more usage
than they report, but they also do not break results down by language or
registry. No prior study focuses on the npm ecosystem specifically, even
though it is the world's largest package registry with more than two million
packages and tens of billions of weekly downloads~\cite{chatzidimitriou2018}.

A second gap is about where in the world developers are using these tools.
Prior work has shown that open-source participation varies significantly across
countries~\cite{wachs2022}, yet no study has compared AI coding agent adoption
between developer communities in the Global North and the Global South.
Egypt is a particularly interesting case. It has one of the fastest-growing
developer communities on GitHub in the Middle East and Africa, but Egyptian
developers face different conditions than developers in, say, the United
States---different infrastructure, different access to paid AI tools, and
different norms around automation. Whether these differences lead to lower AI
adoption---and by how much---is an open question that this thesis addresses.

% ------------------------------------------------------------
\section{Research Questions}
\label{sec:intro:rqs}
% ------------------------------------------------------------

This thesis is structured around three research questions (RQs):

\begin{description}
  \item[RQ1] \textbf{Scope of adoption.}
    What proportion of active JavaScript, TypeScript, and npm repositories on
    GitHub created since February 2025 show evidence of AI coding agent use,
    and how does this vary by package manager ecosystem (npm, Yarn, pnpm, Bun),
    primary language, and repository size?

  \item[RQ2] \textbf{Signal undercounting.}
    To what extent does relying on a single detection signal underestimate AI
    coding agent adoption compared with a multi-signal approach, and which
    signal types contribute most to detection coverage in the JS/TS/npm
    ecosystem?

  \item[RQ3] \textbf{Community comparison.}
    Do repositories with Egyptian-affiliated contributors show significantly
    different AI coding agent adoption rates than the global JS/TS/npm
    baseline, and what factors---tool accessibility, repository size, community
    type---correlate with any observed difference?
\end{description}

These three questions address the scope, the measurement reliability, and the
geographic dimension of AI coding agent adoption---the three gaps identified
in the literature review (Chapter~\ref{chap:litrev}).

% ------------------------------------------------------------
\section{Scope and Delimitations}
\label{sec:intro:scope}
% ------------------------------------------------------------

The study targets \emph{public} GitHub repositories that meet all of the
following criteria: created on or after 24 February 2025 (the start of the
high-growth agentic era); at least four stars (to exclude toy repositories
and forks of popular projects); at least fifty commits (to ensure a meaningful
development history); not themselves forks; and with at least 175\,000 bytes
of combined JavaScript and TypeScript source code (to confirm that JS/TS is
the primary language of the project).

\emph{Ecosystem} is determined by the lockfile present at the repository
root: \texttt{package-lock.json} (npm), \texttt{yarn.lock} (Yarn),
\texttt{pnpm-lock.yaml} (pnpm), \texttt{bun.lockb} (Bun). Repositories with
\texttt{package.json} but no lockfile are classified as \texttt{npm\_no\_lock};
repositories with \texttt{deno.json} are classified as \texttt{deno};
repositories with neither are classified as \texttt{vanilla}. The thesis
focuses on the four lockfile-confirmed ecosystems and \texttt{npm\_no\_lock},
which together represent the main npm dependency surface.

\emph{AI coding agent adoption} is detected using four families of signals
adapted from Robbes et al.~\cite{robbesHeuristics2026} and Galster et al.\
\cite{galsterConfig2026}: (1) configuration files (\texttt{CLAUDE.md},
\texttt{.cursorrules}, \texttt{AGENTS.md}, etc.); (2) commit message patterns
left by agents (\texttt{Co-authored-by: Claude},
\texttt{Generated with Claude Code}, etc.); (3) pull request body patterns;
and (4) AI-specific npm dependency declarations (e.g.\
\texttt{@anthropic-ai/claude-code}, \texttt{@cursor/agent}). Repositories
are classified as \emph{agentic} if they show at least two distinct signal
types, \emph{AI-assisted} if they show exactly one type, and \emph{none}
otherwise.

The Egyptian corpus is built by matching contributor GitHub profiles to
Egypt using location fields, biography text, and organisational affiliations.
The comparison is observational; no causal claims are made.

% ------------------------------------------------------------
\section{Methodology Overview}
\label{sec:intro:method}
% ------------------------------------------------------------

The study follows a seven-stage pipeline:

\begin{enumerate}
  \item \textbf{Global corpus collection} (Stage 1a): collect all JS/TS
    repositories on GitHub that meet the inclusion criteria using the
    GitHub GraphQL Search API with date-windowed pagination to stay within the
    1{,}000-result-per-query limit.

  \item \textbf{Ecosystem categorisation} (Stage 1b): classify each repository
    by its package manager using root-level lockfile detection.

  \item \textbf{Stratified sampling} (Stage 1c): draw a stratified random
    sample proportional to ecosystem × language × size stratum, targeting
    approximately 3{,}500--5{,}000 repositories for Stage 4 analysis.

  \item \textbf{Egyptian corpus construction} (Stage 2): identify repositories
    whose contributors include at least one developer with a confirmed
    Egyptian affiliation, using a nine-signal location inference pipeline.

  \item \textbf{Signal detection} (Stage 4): scan each sampled repository for
    the four families of AI coding agent signals; output per-repository signal
    vectors and adoption-level classifications.

  \item \textbf{Statistical comparison} (Stage 5): compare global and Egyptian
    adoption rates using Fisher's exact test; compare repository characteristics
    (stars, commits, age) using Mann-Whitney U tests; apply Benjamini-Hochberg
    correction for multiple comparisons.

  \item \textbf{Synthesis} (Stage 6): interpret quantitative results in the
    context of the literature; identify qualitative patterns by inspecting
    configuration file content and commit message phrasing.
\end{enumerate}

The pipeline is implemented entirely in Python 3.11, using the GitHub GraphQL
and REST APIs. All scripts, intermediate data files, and the final corpus are
archived in the project repository to ensure reproducibility.

% ------------------------------------------------------------
\section{Contributions}
\label{sec:intro:contributions}
% ------------------------------------------------------------

This thesis makes three contributions to the empirical software engineering
literature:

\begin{enumerate}
  \item \textbf{First ecosystem-scoped adoption estimate.} By focusing the
    corpus on the JS/TS/npm ecosystem and classifying repositories by package
    manager, this study provides adoption estimates that practitioners and
    tool vendors can map directly to a specific technical community, rather
    than averaging across all GitHub languages.

  \item \textbf{Multi-signal undercounting quantification.} By comparing
    single-signal and four-signal detection on the same corpus, the study
    provides a concrete estimate of how much adoption is missed when only
    one detection signal is used---directly addressing the measurement concern
    raised by Khosravani and Mockus~\cite{khosravani2026}.

  \item \textbf{Cross-community comparison.} The Egyptian-versus-global
    comparison provides the first empirical data point on AI coding agent
    adoption in a Middle Eastern/African developer community, contributing
    evidence on equitable access to AI development tools across the Global
    South.
\end{enumerate}

% ------------------------------------------------------------
\section{Thesis Structure}
\label{sec:intro:structure}
% ------------------------------------------------------------

The remainder of this thesis is organised as follows.
Chapter~\ref{chap:litrev} reviews the relevant literature across six themes:
AI coding agent definitions, empirical adoption studies, detection methodology,
developer interaction research, the JS/TS/npm ecosystem as a research context,
and geographic dimensions of open-source development.
Chapter~\ref{chap:methodology} describes the data collection pipeline, corpus
construction, signal detection algorithm, and statistical analysis plan in
full technical detail.
Chapter~\ref{chap:results} presents the quantitative findings for each of the
three research questions.
Chapter~\ref{chap:discussion} interprets the findings, discusses threats to
validity, and considers the implications for tool developers, open-source
maintainers, and researchers studying AI-augmented development.
Chapter~\ref{chap:conclusion} summarises the contributions and outlines
directions for future work.
